\hypertarget{scenario-1-building-a-model-of-cancer-cell-fate-in-response-to-tnf-from-scratch}{%
\section{Scenario 1: Building a model of cancer cell fate in response to
TNF from
scratch}\label{scenario-1-building-a-model-of-cancer-cell-fate-in-response-to-tnf-from-scratch}}

\hypertarget{prompt}{%
\subsection{Prompt:}\label{prompt}}

You are my multiscale modelling assistant. Together we are creating a
mutliscale multicellular model of cancer cell fate in response to TNF.
In order to chieve this we are going to use NeKo, MaBoSS and PhysiCell
to create a simulation where an initial group of cancer cell will grow
while TNF diffuses in the microenvironment. When the cells sense the TNF
they will activate the input nodes of their boolean network. According
to the final state of the network, they will perform behaviors such as
die for apoptosis or proliferte. Those are he steps to follow: 1) create
a network using Omnipath that describes cancer cell fates in response to
tnf. This network should contains TNF sensing receptr as well asgenes
responsible for proliferation and apoptosis. Once create the network we
have to remove the bimodal interactions and check the connectivity. If
the network is connected, export it in bnet file. After that generate
the maboss file (.bnd and .cfg), create a simulation and select some
nodes to identify the output, to represent apoptotic state and
proliferative state. Once set those nodes as output, we run the maboss
simulation, get the results and study the outcome (does the network make
sense? are the states biologically feasible?). This is the first part.
In the second part we will build the PhysiCell simulation and add maboss
to the cancer cell type. Is everything clear?

\hypertarget{final-results}{%
\subsection{Final results:}\label{final-results}}

Here's what I changed and exported.

TNF diffusion from edges

\begin{itemize}
\tightlist
\item
  TNF substrate: diffusion\_coefficient=10000 µm²/min (slower),
  decay\_rate=0.01 1/min, initial=0.
\item
  Edge source: Dirichlet boundary enabled, boundary value=1.0 (TNF
  enters from domain edges and diffuses inward).
\item
  Cancer cells: TNF uptake\_rate=0.01 1/min; no secretion.
\end{itemize}

PhysiBoSS coupling (unchanged)

\begin{itemize}
\tightlist
\item
  Input link: TNF signal → boolean node TNF (threshold=0.1).
\item
  Output links: CASP3 → apoptosis\_rate (base 1e-4, active 0.02); E2F1 →
  proliferation\_rate (base 0.001, active 0.02).
\item
  Intracellular: dt=6 min, scaling=1, inheritance\_global=false.
\end{itemize}

Domain and run

\begin{itemize}
\tightlist
\item
  Size: 1000 x 1000 x 20 µm, dx=20; 2D slab; duration: 24 h.
\item
  Cell type: cancer\_cell (Ki67\_basic) with default volumes/motility;
  PhysiBoSS model attached (output.bnd/output.cfg).
\end{itemize}

Exports and files

\begin{itemize}
\tightlist
\item
  MaBoSS: output.bnd, output.cfg (already in workspace).
\item
  PhysiCell XML:
  \href{vscode-file://vscode-app/usr/share/code/resources/app/out/vs/code/electron-browser/workbench/workbench.html}{PhysiCell\_settings.xml}
  (exported at project root). Use this file with your simulator; if you
  want it in gpt\_5/Scenario\_1, I can place a copy there.
\end{itemize}

Global summary of what we achieved

\begin{itemize}
\tightlist
\item
  Built an OmniPath TNF-fate network, removed bimodal edges, verified
  connectivity, exported .bnet, generated MaBoSS .bnd/.cfg.
\item
  Selected outputs (CASP3=apoptosis, E2F1=proliferation), ran MaBoSS;
  baseline under TNF favored NF-κB survival; mutations (e.g., XIAP OFF)
  modestly increased apoptosis.
\item
  Created a PhysiCell session, added TNF substrate and a cancer cell
  type, integrated the MaBoSS model via PhysiBoSS, wired TNF signal to
  the network and network outputs to cell behaviors.
\item
  Adjusted transport to model slow TNF diffusion from edges toward the
  center; exported an XML ready to run.
\end{itemize}

Want me to:

\begin{itemize}
\tightlist
\item
  Copy the exported XML into gpt\_5/Scenario\_1 and add a basic seeding
  (e.g., 200 cells in a central disk)?
\item
  Add a gentle TNF boundary value (e.g., 0.2 instead of 1.0) or
  time-varying pulse?
\end{itemize}

\hypertarget{scenario-2-building-a-set-of-rules-for-an-existing-physicell-simulation-by-extracting-the-context-directly-from-the-configuration-file}{%
\section{Scenario 2: Building a set of rules for an existing PhysiCell
simulation by extracting the context directly from the configuration
file}\label{scenario-2-building-a-set-of-rules-for-an-existing-physicell-simulation-by-extracting-the-context-directly-from-the-configuration-file}}

\hypertarget{prompt-1}{%
\subsection{Prompt:}\label{prompt-1}}

I have a PhysiCell configuration file that I would like to extend with
biologically accurate rules. Load the PhysiCell configuration file,
extract the biological scenario, extract the substrates, cell types and
custom parameters that compose the simulation and, carefully examining
the available signals and behaviors, add a series of rules to improve
the model.

\hypertarget{final-results-1}{%
\subsection{Final Results:}\label{final-results-1}}

Here's what I did and what I found.

Biological scenario

\begin{itemize}
\tightlist
\item
  Tumor--immune microenvironment with malignant epithelial cells,
  macrophages (M0/M1/M2), and T cells (effector/exhausted).
\item
  Substrates: oxygen (Dirichlet at boundaries), apoptotic/necrotic
  debris (macrophage attractants), cytokines (pro- vs anti-inflammatory
  factors) driving polarization and T-cell function.
\end{itemize}

Extracted components

\begin{itemize}
\tightlist
\item
  Substrates: oxygen, apoptotic debris, necrotic debris,
  pro-inflammatory factor, anti-inflammatory factor.
\item
  Cell types: malignant epithelial cell, M0 macrophage, M1 macrophage,
  M2 macrophage, effector T cell, exhausted T cell.
\item
  Custom parameters:

  \begin{itemize}
  \tightlist
  \item
    \texttt{user\_parameters.number\_of\_cells\ =\ 0}
  \item
    \texttt{custom\_data.sample} per cell type (malignant: 1.0; others:
    0.0)
  \item
    Existing ruleset reference enabled: \texttt{cell\_rules.csv} at
  \end{itemize}
\end{itemize}

\textbf{GitHub Copilot Response:}

Added biologically grounded rules

\begin{itemize}
\tightlist
\item
  Malignant epithelial cell:

  \begin{itemize}
  \tightlist
  \item
    Oxygen increases → cycle entry (half-max \textasciitilde20, Hill 4)
  \item
    Oxygen decreases → necrosis (half-max \textasciitilde5, Hill 6)
  \item
    Pro-inflammatory factor increases → immunogenicity to effector T
  \item
    Anti-inflammatory factor decreases → immunogenicity to effector T
  \end{itemize}
\item
  M0 macrophage:

  \begin{itemize}
  \tightlist
  \item
    Pro-inflammatory factor increases → transition to M1
  \item
    Anti-inflammatory factor increases → transition to M2
  \item
    Apoptotic/necrotic debris increases → phagocytose apoptotic/necrotic
    cells
  \end{itemize}
\item
  M1 macrophage:

  \begin{itemize}
  \tightlist
  \item
    Contact with malignant epithelial cell increases → attack malignant
    epithelial cell
  \item
    Pro-inflammatory factor increases → attack malignant epithelial cell
    and chemotaxis to pro-inflammatory factor
  \item
    Anti-inflammatory factor decreases → attack malignant epithelial
    cell
  \end{itemize}
\item
  M2 macrophage:

  \begin{itemize}
  \tightlist
  \item
    Anti-inflammatory factor increases → anti-inflammatory factor
    secretion (positive feedback)
  \item
    Anti-inflammatory factor decreases → reduce attack on malignant
    epithelial cell
  \item
    Apoptotic debris increases → phagocytose apoptotic cells
  \end{itemize}
\item
  Effector T cell:

  \begin{itemize}
  \tightlist
  \item
    Contact with malignant epithelial cell increases → attack malignant
    epithelial cell
  \item
    Pro-inflammatory factor increases / anti-inflammatory factor
    decreases → attack malignant epithelial cell
  \item
    Pro-inflammatory factor increases → migration speed
  \item
    Anti-inflammatory factor increases → transition to exhausted T cell
  \end{itemize}
\item
  Exhausted T cell:

  \begin{itemize}
  \tightlist
  \item
    Pro-inflammatory factor increases → transition to effector T cell
  \item
    Anti-inflammatory factor decreases → migration speed
  \end{itemize}
\end{itemize}

Exports created (in the workspace root)

\begin{itemize}
\tightlist
\item
  \texttt{cell\_rules\_gpt5\_scenario2.csv} (22 rules)
\item
  \texttt{PhysiCell\_settings\_gpt5\_scenario2\_with\_rules.xml}
  (updated config referencing the new rules)
\end{itemize}

How to use

\begin{itemize}
\tightlist
\item
  Option A: Use
  \texttt{PhysiCell\_settings\_gpt5\_scenario2\_with\_rules.xml}
  directly.
\item
  Option B: Keep your original XML and point its
  \texttt{\textless{}cell\_rules\textgreater{}\textless{}ruleset\textgreater{}}
  to \texttt{cell\_rules\_gpt5\_scenario2.csv} (and ensure the CSV is
  placed where the XML expects it), then run your simulation.
\end{itemize}

\hypertarget{scenario-3-iteratevely-refine-a-maboss-model-through-neko-using-the-basic-knowledge-of-the-llm}{%
\section{Scenario 3: iteratevely refine a MaBoSS model through NeKo
using the basic knowledge of the
LLM}\label{scenario-3-iteratevely-refine-a-maboss-model-through-neko-using-the-basic-knowledge-of-the-llm}}

\hypertarget{prompt-2}{%
\subsection{Prompt:}\label{prompt-2}}

I want to create a Boolean model that describes how TNF affects cancer
cell fate. In order to do so, I want you to create a network using
Omnipath as interactions database. Once created the network, remove the
bimodal interactions, export the network in bnet format and convert it
into a MaBoSS model. Run the MaBoSS model (make sure to set thread count
to 10 and select a subset of output nodes to reduce computation time)
and get the results. If the results are biologically accurate and
relevant, test some known mutations end verify the accuracy of the
results. If the network is not satisfying, check important pathways
within the network and individuate nodes or interactions that may not be
relevant and bring inaccuracies in the model. Finally, make sure the
final network always contains relevant nodes, important for our
scenario. Iterate the process until the MaBoSS model is accurate enough.

\hypertarget{final-results-2}{%
\subsection{Final results:}\label{final-results-2}}

Here's a concise summary of what was done, how the curated model
differs, and why the logic changes are biologically supported.

What was achieved

\begin{itemize}
\tightlist
\item
  Built a TNF-driven Boolean network from Omnipath, removed bimodal
  edges, exported to BNET, converted to MaBoSS.
\item
  Iteratively curated key logical rules for NF-κB, apoptosis, and
  necroptosis.
\item
  Set TNF stimulus, selected outputs {[}RELA, CASP3, MLKL{]}, ran MaBoSS
  with thread\_count=10.
\item
  Achieved baseline consistent with TNF-induced NF-κB activation
  (RELA=1, CASP3=0, MLKL=0).
\item
  Validated NF-κB loss with IKBKB OFF; explored necro/apoptosis toggling
  via cFLIP/IAP logic.
\end{itemize}

Key differences: original vs curated (selected rules)

\begin{itemize}
\tightlist
\item
  TNF

  \begin{itemize}
  \tightlist
  \item
    Original: TNF = (TNFAIP3) \& !(CFLAR \textbar{} PTEN)
  \item
    Curated: TNF = TNF (hold stimulus constant)
  \item
    Rationale: treat TNF as external input during stimulation.
  \end{itemize}
\item
  TRADD

  \begin{itemize}
  \tightlist
  \item
    Original: TRADD = TNFRSF1A \& !(TNF \textbar{} CFLAR)
  \item
    Curated: TRADD = TNFRSF1A \& !CFLAR
  \item
    Rationale: TRADD should be recruitable under TNF; cFLIP can inhibit
    DISC assembly.
  \end{itemize}
\item
  CASP8 (DISC)

  \begin{itemize}
  \tightlist
  \item
    Original: CASP8 = (TNF \textbar{} TNFRSF1A \textbar{} TRADD
    \textbar{} TRAF2 \textbar{} FADD) \& !(CFLAR \textbar{} BIRC2
    \textbar{} BIRC3 \textbar{} NFKB1 \textbar{} RELA)
  \item
    Curated: CASP8 = (TNFRSF1A \& TRADD \& FADD) \& !(CFLAR \textbar{}
    BIRC2 \textbar{} BIRC3)
  \item
    Rationale: canonical DISC requires TNFRSF1A/TRADD/FADD; cFLIP and
    cIAPs suppress caspase-8.
  \end{itemize}
\item
  IKK/NF-κB activation and IκBα

  \begin{itemize}
  \tightlist
  \item
    RELA

    \begin{itemize}
    \tightlist
    \item
      Original: required IKK plus upstream nodes and !CASP3 or !NFKBIA
      variants
    \item
      Curated: RELA = (IKBKB \& IKBKG \& CHUK) \& (TNF \textbar{}
      TNFRSF1A \textbar{} TRADD \textbar{} TRAF2 \textbar{} RIPK1)
    \end{itemize}
  \item
    NFKBIA

    \begin{itemize}
    \tightlist
    \item
      Original: NFKBIA = (TRAF2 \textbar{} CASP3 \textbar{} IKBKG
      \textbar{} NFKB1) \& !(TNF \textbar{} TNFRSF1A \textbar{} RELA
      \textbar{} TNFAIP3)
    \item
      Curated: NFKBIA = (RELA \textbar{} NFKB1) \& !(IKBKB \& IKBKG \&
      CHUK)
    \end{itemize}
  \item
    Rationale: IKK complex (IKKβ/IKBKB, NEMO/IKBKG, IKKα/CHUK) activates
    NF-κB; IκBα is an NF-κB target and is degraded by IKK.
  \end{itemize}
\item
  A20 (TNFAIP3)

  \begin{itemize}
  \tightlist
  \item
    Original: TNFAIP3 = (TNF \textbar{} TNFRSF1A \textbar{} NFKB1
    \textbar{} RELA)
  \item
    Curated: TNFAIP3 = (NFKB1 \textbar{} RELA)
  \item
    Rationale: A20 is an NF-κB--induced negative feedback, not a direct
    TNF output.
  \end{itemize}
\item
  TRAF2 and IKBKG (NEMO)

  \begin{itemize}
  \tightlist
  \item
    Original: TRAF2 inhibited by TNFAIP3; IKBKG inhibited by TNFAIP3
  \item
    Curated: TRAF2 = (TNFRSF1A \textbar{} RELA); IKBKG = (TNF \textbar{}
    TNFRSF1A \textbar{} TRADD \textbar{} TRAF2 \textbar{} RIPK1
    \textbar{} FADD \textbar{} CASP8)
  \item
    Rationale: relieve overly strong A20 gating to allow initial IKK
    activation under TNF.
  \end{itemize}
\item
  cFLIP (CFLAR)

  \begin{itemize}
  \tightlist
  \item
    Original: constitutively ON via negations (blocking DISC)
  \item
    Curated: CFLAR = (NFKB1 \textbar{} RELA \textbar{} AKT1)
  \item
    Rationale: cFLIP is NF-κB/AKT-inducible; initial DISC can form
    before feedback.
  \end{itemize}
\item
  cIAPs (BIRC2/3)

  \begin{itemize}
  \tightlist
  \item
    Original: BIRC2/3 activated broadly and also drove RIPK1
  \item
    Curated: BIRC2 = (NFKB1 \textbar{} RELA), BIRC3 = (NFKB1 \textbar{}
    RELA)
  \item
    Rationale: cIAPs are NF-κB targets; they stabilize complex I and
    suppress necrosome.
  \end{itemize}
\item
  RIPK1 and RIPK3/MLKL (necroptosis)

  \begin{itemize}
  \tightlist
  \item
    RIPK1

    \begin{itemize}
    \tightlist
    \item
      Original: (TNFRSF1A \textbar{} TRADD \textbar{} TRAF2 \textbar{}
      FADD \textbar{} BIRC2 \textbar{} BIRC3) \& !(CFLAR \textbar{}
      CASP8 \textbar{} TNFAIP3)
    \item
      Curated: (TNFRSF1A \textbar{} TRADD \textbar{} TRAF2 \textbar{}
      FADD) \& !(CFLAR \& CASP8)
    \end{itemize}
  \item
    RIPK3

    \begin{itemize}
    \tightlist
    \item
      Original: (MAP3K7) \& !CASP8
    \item
      Curated: RIPK3 = RIPK1 \& !(CASP8 \textbar{} BIRC2 \textbar{}
      BIRC3)
    \end{itemize}
  \item
    MLKL unchanged: MLKL = RIPK3
  \item
    Rationale: necroptosis requires RIPK1→RIPK3 when caspase-8 is
    inactive; cIAPs block necrosome formation; TAK1 alone should not
    drive RIPK3.
  \end{itemize}
\item
  Apoptosome and executioner caspases

  \begin{itemize}
  \tightlist
  \item
    CASP3

    \begin{itemize}
    \tightlist
    \item
      Original: mixed sources with strong inhibitors
    \item
      Curated: CASP3 = (CASP8 \textbar{} (APAF1 \& CYCS)) \& !XIAP
    \end{itemize}
  \item
    APAF1

    \begin{itemize}
    \tightlist
    \item
      Curated: APAF1 = (BAX \textbar{} CYCS) \& !XIAP
    \end{itemize}
  \item
    CASP9

    \begin{itemize}
    \tightlist
    \item
      Curated: CASP9 = (APAF1 \textbar{} CYCS) \& !BCL2L1
    \end{itemize}
  \item
    Rationale: align with intrinsic apoptosis control by XIAP and BCL2
    family, and extrinsic via CASP8.
  \end{itemize}
\end{itemize}

Behavior after curation (representative)

\begin{itemize}
\tightlist
\item
  Baseline TNF ON: RELA=1 (NF-κB active), CASP3=0, MLKL=0.
\item
  IKBKB OFF: NF-κB off (\texttt{\textless{}nil\textgreater{}}=1) as
  expected.
\item
  RIPK3 OFF: NF-κB remains on, necro off (expected).
\item
  CASP8 OFF: necro not engaged unless IAP/cFLIP axis is also low
  (biologically consistent with requirement for cIAP loss or cFLIP
  down).
\end{itemize}

References supporting the curated logic

\begin{itemize}
\tightlist
\item
  NF-κB activation via IKK complex; IκBα negative feedback and
  IKK-mediated degradation:

  \begin{itemize}
  \tightlist
  \item
    Hayden \& Ghosh. NF-κB in immunobiology. Cell Res 2011; Karin \&
    Ben-Neriah. NF-κB at 20. Nat Immunol 2000; Hoffmann et al.~The
    IκB--NF-κB signaling module. Science 2002.
  \end{itemize}
\item
  A20 (TNFAIP3) as NF-κB--induced negative regulator of TNFR signaling
  and RIPK1 ubiquitination:

  \begin{itemize}
  \tightlist
  \item
    Wertz et al.~De-ubiquitination and A20 function. Nature 2004; Boone
    et al.~A20 and NF-κB. Nat Immunol 2004; Shembade \& Harhaj. A20 in
    NF-κB termination. Trends Immunol 2012.
  \end{itemize}
\item
  cFLIP as NF-κB/AKT-inducible DISC inhibitor:

  \begin{itemize}
  \tightlist
  \item
    Krueger et al.~FLIP inhibits caspase-8. Curr Biol 2001; Micheau et
    al.~DISC regulation. J Biol Chem 2002; Kataoka et al.~FLIP
    expression and NF-κB. Mol Cell 2000.
  \end{itemize}
\item
  cIAP1/2 (BIRC2/3) promote complex I (NF-κB) and suppress
  RIPK1-dependent cell death:

  \begin{itemize}
  \tightlist
  \item
    Varfolomeev \& Vucic. cIAPs in TNF signaling. Cell Death Differ
    2011; Bertrand et al.~cIAPs and RIPK1. Cell 2008.
  \end{itemize}
\item
  Caspase-8 suppresses necroptosis; RIPK1→RIPK3→MLKL drives necroptosis
  when caspase-8 is inhibited:

  \begin{itemize}
  \tightlist
  \item
    Oberst et al.~Caspase-8 and necroptosis. Nature 2011; He et al.~MLKL
    mediates necroptosis. Cell 2009/2013; Sun et al.~MLKL requirement.
    Cell 2012; Degterev et al.~Necrostatin. Nat Chem Biol 2005.
  \end{itemize}
\item
  Extrinsic and intrinsic apoptosis wiring:

  \begin{itemize}
  \tightlist
  \item
    Ashkenazi \& Dixit. Death receptors. Science 1998; Li et al.~Apaf-1
    and cytochrome c.~Cell 1997; Zou et al.~Apaf-1. Cell 1997; Deveraux
    \& Reed. IAPs inhibit caspases. Genes Dev 1999; Youle \& Strasser.
    BCL-2 family. Nat Rev Mol Cell Biol 2008.
  \end{itemize}
\end{itemize}

If you want, I can package the curated BNET and the MaBoSS outputs, or
continue refining to make necro/apoptosis engage under appropriate
combined perturbations.
